% v19 G5_7 obstruction. Dependencies: mh18:matching, tl15:one-high-full,
% ml12:dichotomy, ls11:low-puncture, td12:low-entry.
% Exact arithmetic certificate: verify_g5_7_obstruction_v19.py.
\section{三角形を共有する菱形に一葉を付けた支持の排除}
\label{g719:section}

高支持を
\[
 H=\{01,02,12,13,23,04\}
\]
とする．これは標準代表 $G5\_7=\{01,02,03,12,13,24\}$ を
旧頂点 $(2,0,1,3,4)$ の順に新頂点 $(0,1,2,3,4)$ へ写したものである．
以下では葉 $4$ の標準形の低 rays を $B,C,D$ と呼び，$A=0$ とする．
三つの低 rays のうち $0$ との辺が低分母であるものを $m$ と書く．
従って $m$ は物理的には頂点 $3$ であるが，標準形内での位置は
$B,C,D$ の三通りを全て保つ．残る二頂点は互換である．

\begin{lemma}[二高 core の二つの小さい正 reverse は共存しない]
\label{g719:small-positive}
正の整数二高局所模型の二つの低行を取り，各低行中の唯一の正 entry を
$e_1,e_2$ とする．$e_1,e_2\le2$ は不可能である．
\end{lemma}
\begin{proof}
二高標準形ではこれらは $r,T$ であり，低行は
$(r,-r-1,0)$ と $(-S,T,S-T-1)$ である．
正の枝では $S>r$ および $rT>(r+1)S$ なので $T>r+1\ge2$ となる．
負の枝では $S=T<r$ なので，$r,T\le2$ のもとでは
$r=2,S=T=1$ に限られる．この場合 $a=s+1,b=z+1$，
$x=s/2,y=z/3$，$A=1+s/2,B=1+5z/3$ であり，
$AB=ab-1$ は
\[
 (s-4)(z+3)=-6
\]
と同値である．整数性から $s\ge2$ は偶数，$z\ge3$ は三の倍数である．
$s=2$ なら左辺は $-12$ 以下，$s\ge4$ なら非負であり矛盾する．
負の枝を捨てず，その唯一の小さい場合をここで明示的に排除した．
\end{proof}

葉の完全標準形を
\[
 V=P^4=\begin{pmatrix}
 0&1-Q&-h-1&h\\
 -1&0&0&0\\
 -u&u-1-c&0&c\\
 v&-v-1-d&d&0
 \end{pmatrix},
\qquad Q=th(h+1)+1,
\]
\[
 u=th,\quad v=t(h+1),\quad cd=0
\]
と書く．$h,t\ge1$ であり，行列の順序は $(0,B,C,D)$ である．
$N_{ij}=n_{i4;j}$ とし，各行の三成分は非負で総和五である．
\[
 \ell_0=Q,\quad \ell_B=\ell_C=\ell_D=1,\qquad
 R_{ij}=p^i_{4j}=\ell_i(1-N_{ij})-V_{ij},\quad
 X_{ij}=p^i_{j4}
\]
と置けば
\begin{equation}\label{g719:basic}
 R_{ij}X_{ij}=\ell_iq_{ij}-1,
 \qquad
 \sum_{j\ne i}X_{ij}+\sum_{j\ne i}V_{ji}=-2\ell_i.
\end{equation}
第二式は辺 $i4$ の $D$ 相殺である．特に
\begin{equation}\label{g719:hub-D}
 \sum_{j=B,C,D}X_{0j}=1-2Q-t.
\end{equation}
中央 core $0$ は三高なので，matching 定理から行 $R_0$ はちょうど一つの
正 entry を持つ．core $m$ は二高で，その低 rays は $0,4$ である．
従って
\begin{equation}\label{g719:low-pair}
 R_{m0}\le0,\quad X_{m0}\le0,\quad R_{m0}X_{m0}=0,\quad
 R_{0m}X_{0m}=Q-1,\quad X_{0m}+X_{m0}\in\{-4,-3,-2,-1,0,1\}.
\end{equation}

\begin{lemma}[正の最大 divisor の二枝の排除]
\label{g719:positive-branches}
次の二つの枝はいずれも不可能である：
\[
 (m=B,\ N_{0B}=1),\qquad
 (m=D,\ h=1,\ N_{0D}=0).
\]
\end{lemma}
\begin{proof}
両方の場合で $R_{0m}=Q-1>0$，$X_{0m}=1$ である．
中央の matching と低行 $m$ の和 $-1$ から，残る二高 rays を $a,b$ と書くと
\[
 p^0_{ma}=p^0_{mb}=-1.
\]
従って低辺 $0m$ の clutch により
$p^m_{0a}=2-n_{0m;a}\le2$，$p^m_{0b}=2-n_{0m;b}\le2$ である．

$m=B$ の場合，葉の $B$ 行は $(-1,0,0)$ である．
従って core $m$ の低行 $4$ の二つの高 entries は
$1-N_{ma},1-N_{mb}$ であり，唯一の正 entry は一である．
低行 $0$ の正 entry は二以下なので，補題~\ref{g719:small-positive} に反する．

$m=D,h=1$ の場合には残る高 rays を葉の $B,C$ の順に置き，
$A=q_{0B},B'=q_{0C}$ とする．中央の正高間 entries を
$H=p^0_{BC}>0,K=p^0_{CB}>0$ と書けば，row sums と pair equations は
\[
 p^0_{B4}=-H-1,\quad p^0_{C4}=-K-1,\quad HK=AB'-1.
\]
式~\eqref{g719:hub-D} と $t=(Q-1)/2$ から
\[
 H+K=\frac{5(Q-1)}2.
\]
$R_{0B},R_{0C}<0$ なので $N_{0B},N_{0C}\ge2$ であり，
$N_{0B}+N_{0C}=5$ から分配は $(2,3)$ または $(3,2)$ である．
前者では $R_{0B}=-1$，従って $H=QA-2\ge3Q-2>5(Q-1)/2$ となり矛盾する．
ここで高分母は奇数なので $A\ge3$ を用いた．
後者では
\[
 R_{0B}=-(Q+1),\quad R_{0C}=-(Q-2),\quad
 H+1=\frac{QA-1}{Q+1}.
\]
従って $A=k(Q+1)-1$，$H=Qk-2$ であり，$H,K>0$ により $k=1,2$ に限る．
それぞれについて
\[
\begin{array}{c|cccc|c}
 k&A&B'&H&K&HK-(AB'-1)\\\hline
 1&Q&(3Q-5)/2&Q-2&(3Q-1)/2&2-Q\\
 2&2Q+1&(Q-1)/2&2Q-2&(Q-1)/2&(5-3Q)/2
\end{array}
\]
となり，$Q\ge3$ では最後の列が零にならない．
\end{proof}

\begin{proposition}[葉分母と配置の完全な有限帰着]
\label{g719:leaf-bound}
上記支持で正の整数局所模型，marked clutch，非負 node 配分（各辺総和五），
および全ての対向 $D$ 相殺が成立するなら $Q\le21$ であり，
さらに次の表のいずれかに属する：
\[
\begin{array}{c|c|c|c}
 m&N_{0m}&R_{0m}&\text{追加条件}\\\hline
 B&2&-1&Q\le5\\
 C&1&h+1&th\le4\\
 C&2&-1&h=t=1\\
 D&1&-h&t(h+1)\le4
\end{array}
\]
この帰着では他の高分母，および $c,d$ に上限を置かない．
\end{proposition}
\begin{proof}
式~\eqref{g719:low-pair} により $R_{0m}$ は $Q-1$ の非零 signed divisor である．
$m=B$ なら $R_{0B}=Q(2-N_{0B})-1$ なので，その絶対値から
$N_{0B}=1,2$ に限る．前者は補題~\ref{g719:positive-branches} で排除済み．
後者では $X_{0B}=-(Q-1)$ であり，$X_{B0}\le0$ と clutch から
$-(Q-1)\ge-4$，従って $Q\le5$ である．

$m=C$ なら $R_{0C}=Q(1-N_{0C})+h+1$ である．
$N_{0C}=0$ または $N_{0C}\ge3$ では絶対値が $Q-1$ より大きい．
$N_{0C}=2$ では $Q-h-1$ が $Q-1$ を割るので $Q-h-1\mid h$，
従って $th(h+1)\le2h$ であり，$h=t=1$ に限る．
$N_{0C}=1$ なら $R_{C0}=u+1-N_{C0}\le0$ と $N_{C0}\le5$ から
$u=th\le4$ である．$h\le u$ より
$Q=u(h+1)+1\le u(u+1)+1\le21$ となる．

$m=D$ なら $R_{0D}=Q(1-N_{0D})-h$ である．$N_{0D}\ge2$ は大きすぎる．
$N_{0D}=0$ なら $Q-h\mid Q-1$，従って $Q-h\mid h-1$ である．
$h\ge2$ では $Q-h>h-1$ なので不可能であり，$h=1$ は
補題~\ref{g719:positive-branches} で排除済みである．
$N_{0D}=1$ なら $X_{0D}=-v$，かつ
$R_{D0}=1-N_{D0}-v<0$ なので $X_{D0}=0$ である．
従って $-v\in\{-4,-3,-2,-1,0,1\}$ より $v\le4$ であり，
$Q=hv+1\le(v-1)v+1\le13$ である．
\end{proof}

\begin{lemma}[全整数の自由係数を保つ有限型検査]
\label{g719:finite-types}
命題~\ref{g719:leaf-bound} の下で，必要な node，matching，共通低 puncture 条件，
および辺 $04$ の $D$ 相殺を全て検査すると，候補は残らない．
自由係数 $c,d$ および $q_{0j}>1$ に人為的な上限を置く必要はない．
\end{lemma}
\begin{proof}
標準形内の $B,C,D$ を以下では $1,2,3$ と書く．
$Q\le21$ より $h(h+1)\le20$ なので $1\le h\le4$，
$1\le t\le\lfloor20/(h(h+1))\rfloor$ である．
$(h,t)$ は十五通りであり，$m=B,C,D$ と合わせて四十五通りとなる．
このうち命題~\ref{g719:leaf-bound} の配置表だけを保つ．
各 $N_i$ は非負三成分で和五なので二十一通りである．
従って式~\eqref{g719:basic} の $R$ は，$c=0$ または $d=0$ の二枝それぞれにおいて，
一つの行の二成分だけが自由整数 $z$ の一次式となる．

$i\ne m$ の中央高辺に対して $R=R_{0i},r=R_{i0}$ と置くと，
同じ marked puncture $4$ における clutch は
\begin{equation}\label{g719:central-equation}
 (Qr+R-\kappa Rr)q_{0i}=R+r,\qquad -4\le\kappa\le1.
\end{equation}
左辺係数が非零なら $q_{0i}$ はただ一つの有理数として決まり，
奇数整数 $\ge3$ と二つの pair 整除を満たすものだけを保つ．
左辺係数が零なら，右辺が零である場合のみ可能であり，
\[
 r=-R,\quad Q-1=-\kappa r,\quad
 q_{0i}=1+\operatorname{lcm}(2,|r|)a,\qquad a\in\ZZ_{>0}
\]
という完全な無限等差数列を保つ．
$i=m$ の低辺では $q_{0m}=1$ と式~\eqref{g719:low-pair} をそのまま用いる．

$1,2,3$ 間の各辺は高分母であり，共通 puncture $4$ は両端で低分母である．
従って補題~\ref{ls11:low-puncture} による reverse partner の必要条件は
\begin{equation}\label{g719:partner}
 r=-s\ne0\quad\text{または}\quad
 rs>0,\quad\min(|r|,|s|)=1,\quad\max(|r|,|s|)\in2\ZZ_{>0}.
\end{equation}
固定 $r\ne0$ に対する $s$ の集合は，一点 $-r$ に加えて，
$|r|$ が偶数なら一点 $\operatorname{sgn}(r)$，
$|r|=1$ なら全半直線 $2a\operatorname{sgn}(r)$，$a\ge1$ である．
従って $s=\pm z+b$ に対しても，一点または偶奇を指定した全半直線に厳密に移される．
変数行の二つの条件の共通部分は，端点と偶奇の比較だけで空・非空を決定できる．
各行の正 entry がちょうど一つという条件も，これらの型の符号で決定できる．
二高 core $m$ では $R_{m0}=0$ も保ち，他の二高 entries は非零で逆符号とする．

この有限整数演算は，同梱の標準 Python ライブラリのみの検算
\nolinkurl{verify_g5_7_obstruction_v19.py} に完全に記載した．
固定行間の条件を通る node 配置は $529$，自由行の partner 型対は $1988$ であり，
そのうち符号適合は $910$，空交差は $750$，非空の型は $160$ となる．
中央辺の $\kappa$ の選択まで展開すると $212$ 型であり，
$D$ 相殺に必要な情報でまとめると次の十三行だけが残る．
表中 $z$ は高分母なので奇数整数 $\ge3$ でなければならず，
$\Delta=\sum_jX_{0j}-(1-2Q-t)$ とする．
\begin{center}
\small
\begin{tabular}{c|c|c|c|c|c|r}
$m$&$(h,t)$&$Q$&$(R_{0B},R_{0C},R_{0D})$&$(q_{0B},q_{0C},q_{0D})$&$D$ 条件の障害&型数\\\hline
$B$&$(1,1)$&3&$(-1,2,-4)$&$(1,z,7)$&$z=1$&12\\
$B$&$(1,1)$&3&$(-1,2,-4)$&$(1,z,3)$&$z=-1$&2\\
$B$&$(1,1)$&3&$(-1,-4,2)$&$(1,7,z)$&$z=1$&6\\
$B$&$(1,1)$&3&$(-1,-4,2)$&$(1,3,z)$&$z=-1$&2\\
$B$&$(1,2)$&5&$(-1,-8,4)$&$(1,5,z)$&$z=-3$&1\\
$C$&$(1,1)$&3&$(-1,2,-4)$&$(z,1,7)$&$z=1$&24\\
$C$&$(1,2)$&5&$(-1,2,-6)$&$(z,1,11)$&$z=1$&24\\
$D$&$(1,1)$&3&$(5,-7,-1)$&$(7,5,1)$&$\Delta=6$&42\\
$D$&$(1,1)$&3&$(-7,5,-1)$&$(5,7,1)$&$\Delta=6$&42\\
$D$&$(1,2)$&5&$(-11,7,-1)$&$(9,3,1)$&$\Delta=5$&12\\
$D$&$(2,1)$&7&$(-15,10,-2)$&$(13,13,1)$&$\Delta=14$&24\\
$D$&$(2,1)$&7&$(-15,10,-2)$&$(13,3,1)$&$\Delta=7$&9\\
$D$&$(3,1)$&13&$(-27,17,-3)$&$(25,21,1)$&$\Delta=26$&12
\end{tabular}
\end{center}
最後の計算では $X_{0j}=(Qq_{0j}-1)/R_{0j}$ を用いるだけである．
全ての無限等差数列では $D$ 方程式の一意的な解が三未満であり，
分母が全て決まった行では $\Delta\ne0$ である．
これで全てが矛盾する．なお型数は branch と node 配分を保持した個数であり，
$c=d=0$ は二枝に現れるので同型類の個数ではない．
\end{proof}

\begin{theorem}[$G5\_7$ の無上界排除]
\label{g719:no-G5-7}
五次 all-plane SNC，標準四直線 cones，固定された marked punctures の枠内で，
局所正 weights，非負 node 配分（各辺総和五），marked clutch，
および対向 $D$ 相殺を全て要求するとき，高分母支持 $G5\_7$ は存在しない．
この結論は分母または numerator の探索 cutoff に依存しない．
\end{theorem}
\begin{proof}
葉を頂点 $4$，その高隣接を $0$ とする．中央 $0$ は三高，
残る非隣接 $m$ は二高，残る二つの cores は三高である．
一高完全標準形の全標識を $m=B,C,D$ に分ければ漏れはない．
命題~\ref{g719:leaf-bound} と補題~\ref{g719:finite-types} が全場合を排除する．
\end{proof}
